%! Rates of Change in Polar Functions — Unit 3, Lesson 3.15 — Guided Notes
\documentclass[10pt]{article}
\usepackage{precalculus-article}
\usepackage{precalculus-boxes}
\newcommand{\TallMath}[1]{$\displaystyle #1\rule[-1.4em]{0pt}{3.2em}$}
\newcommand{\CourseName}{AP Precalculus}
\newcommand{\SchoolYear}{2026--2027}
\newcommand{\MeetingLength}{55 minutes}
\newcommand{\UnitNumberName}{Unit 3: Trigonometric and Polar Functions \quad}
\newcommand{\LessonNumberName}{Lesson 3.15: Rates of Change in Polar Functions}

\begin{document}

\pageheader{Unit 3, Lesson 3.15}{Guided Notes}
\namedateperiod

\vspace{0.10in}

%----------------------------------------------------------------------
% Objectives
%----------------------------------------------------------------------
\begin{objectivebox}
\textbf{I will be able to\ldots}
\begin{enumerate}[leftmargin=1.5em, itemsep=4pt]
  \item \textbf{LO 3.15.A} --- Compute the average rate of change of a polar
    function $r = f(\theta)$ over an interval and interpret it as how quickly
    the radius changes per radian. \writeline
  \item \textbf{LO 3.15.A} --- Classify intervals of a polar function as positive
    and increasing, positive and decreasing, negative and increasing, or negative
    and decreasing, and describe the corresponding curve behavior. \writeline
\end{enumerate}
\end{objectivebox}

\vspace{0.08in}

%----------------------------------------------------------------------
% Vocabulary
%----------------------------------------------------------------------
\begin{vocabbox}
\begin{description}[leftmargin=0pt, itemsep=6pt]
  \item[\termblanklong{rate of change of a polar function}]
    For $r = f(\theta)$, the average rate of change over $[\theta_1, \theta_2]$ is
    $\dfrac{f(\theta_2) - f(\theta_1)}{\theta_2 - \theta_1}$. Measures how quickly
    $r$ changes per \underline{\hspace{2cm}} of angle change.

  \item[\termblanklong{positive and increasing}]
    $r > 0$ and $r$ is \underline{\hspace{2cm}}. The curve moves
    \underline{\hspace{3cm}} from the origin as $\theta$ increases.

  \item[\termblanklong{positive and decreasing}]
    $r > 0$ and $r$ is \underline{\hspace{2cm}}. The curve moves
    \underline{\hspace{3cm}} the origin as $\theta$ increases.

  \item[\termblanklong{negative and increasing}]
    $r < 0$ and $r$ is becoming \underline{\hspace{2.5cm}} negative (closer to zero).
    The curve is traced on the \underline{\hspace{3.5cm}} ray, moving toward the origin.

  \item[\termblanklong{negative and decreasing}]
    $r < 0$ and $r$ is becoming \underline{\hspace{2.5cm}} negative (farther from zero).
    The curve is traced on the opposite ray, moving \underline{\hspace{2.5cm}} the origin.
\end{description}
\end{vocabbox}

\vspace{0.08in}

%----------------------------------------------------------------------
% Hook
%----------------------------------------------------------------------
\begin{hookbox}
\textbf{Entry Question:} The table below shows values of $r = 1 + \sin\theta$
at eight key angles.

\vspace{4pt}
\renewcommand{\arraystretch}{1.4}
\begin{tabular}{|c|c|c|c|c|c|c|c|c|}
\hline
$\theta$ & $0$ & $\pi/4$ & $\pi/2$ & $3\pi/4$ & $\pi$ & $5\pi/4$ & $3\pi/2$ & $7\pi/4$ \\
\hline
$r$ & 1 & $\approx 1.71$ & 2 & $\approx 1.71$ & 1 & $\approx 0.29$ & 0 & $\approx 0.29$ \\
\hline
\end{tabular}

\vspace{8pt}
\textbf{Q1.} On which interval(s) is $r$ \textbf{increasing} (values going up)?

\vspace{4pt}
\underline{\hspace{8cm}}

\vspace{6pt}
\textbf{Q2.} On which interval(s) is $r$ \textbf{decreasing} (values going down)?

\vspace{4pt}
\underline{\hspace{8cm}}

\vspace{6pt}
\textbf{Q3.} What is happening at $\theta = 3\pi/2$, where $r = 0$?

\vspace{4pt}
\underline{\hspace{8cm}}

\end{hookbox}

\vspace{0.08in}

%----------------------------------------------------------------------
% LO 3.15.A — Average Rate of Change
%----------------------------------------------------------------------
\begin{notesbox}{LO 3.15.A --- Average Rate of Change of $r$}

\textbf{Formula:}
$$\text{Average Rate of Change} = \frac{f(\theta_2) - f(\theta_1)}{\theta_2 - \theta_1}$$

This tells us: for every additional \underline{\hspace{1.5cm}} that $\theta$ increases,
$r$ changes by approximately \underline{\hspace{2cm}}.

\vspace{8pt}
\textbf{Worked Example:} $r = 2\sin\theta$

\vspace{4pt}
\renewcommand{\arraystretch}{1.5}
\begin{tabular}{|c|c|c|}
\hline
$\theta$ & $0$ & $\pi/2$ \\
\hline
$r = 2\sin\theta$ & $0$ & $2$ \\
\hline
\end{tabular}

\vspace{8pt}
Average rate of change from $\theta = 0$ to $\theta = \pi/2$:
$$\frac{2\sin(\pi/2) - 2\sin(0)}{\pi/2 - 0} = \frac{2 - 0}{\pi/2} = \frac{4}{\pi}$$

Result $= $ \underline{\hspace{3cm}}

\vspace{6pt}
Interpret: As $\theta$ increases from $0$ to $\pi/2$, $r$ increases by approximately
\underline{\hspace{1.8cm}} per radian on average.

\vspace{8pt}
\textbf{Your turn:} Compute the average rate of change from $\theta = \pi/2$ to $\theta = \pi$
for the same function $r = 2\sin\theta$. ($r(\pi/2) = 2$, $r(\pi) = 0$.)

\vspace{6pt}
ARC $= \dfrac{\rule{1.5cm}{0.4pt} - \rule{1.5cm}{0.4pt}}{\rule{2cm}{0.4pt}} = $ \underline{\hspace{3cm}}

\vspace{4pt}
Is $r$ increasing or decreasing on $[\pi/2, \pi]$? \underline{\hspace{3cm}}

\end{notesbox}

\vspace{0.08in}

%----------------------------------------------------------------------
% LO 3.15.A — Sign and Direction
%----------------------------------------------------------------------
\begin{notesbox}{LO 3.15.A --- Sign and Direction of $r$}

\textbf{Four-case classification.} Fill in the ``What happens in the plane'' column.

\vspace{6pt}
\renewcommand{\arraystretch}{1.7}
\begin{tabular}{|l|l|p{4.5cm}|}
\hline
\textbf{Sign of $r$} & \textbf{Behavior of $r$} & \textbf{What happens in the plane} \\
\hline
$r > 0$ & increasing & \underline{\hspace{4cm}} \\
\hline
$r > 0$ & decreasing & \underline{\hspace{4cm}} \\
\hline
$r < 0$ & increasing (less negative) & \underline{\hspace{4cm}} \\
\hline
$r < 0$ & decreasing (more negative) & \underline{\hspace{4cm}} \\
\hline
\end{tabular}

\vspace{8pt}
\textbf{Key reminder:} When $r < 0$, the point is plotted on the
\underline{\hspace{3.5cm}} direction from the terminal ray, at distance
$|r|$ from the origin.

\end{notesbox}

\vspace{0.08in}

%----------------------------------------------------------------------
% LO 3.15.A — Intervals from a Table
%----------------------------------------------------------------------
\begin{notesbox}{LO 3.15.A --- Reading Intervals from a Table}

\textbf{Example:} The table below shows values of $r = 1 + \cos\theta$ at eight key angles.
For each pair of adjacent columns, identify whether $r$ is positive and increasing (P\&I),
positive and decreasing (P\&D), zero, or negative (N).

\vspace{6pt}
\renewcommand{\arraystretch}{1.4}
\begin{tabular}{|c|c|c|c|c|c|c|c|c|}
\hline
$\theta$ & $0$ & $\pi/4$ & $\pi/2$ & $3\pi/4$ & $\pi$ & $5\pi/4$ & $3\pi/2$ & $7\pi/4$ \\
\hline
$r$ & $2$ & $\approx 1.71$ & $1$ & $\approx 0.29$ & $0$ & $\approx 0.29$ & $1$ & $\approx 1.71$ \\
\hline
\end{tabular}

\vspace{8pt}
\begin{tabular}{|c|p{4cm}|}
\hline
\textbf{Interval} & \textbf{Classification} \\
\hline
$[0, \pi/4]$ & \underline{\hspace{3.5cm}} \\
\hline
$[\pi/4, \pi/2]$ & \underline{\hspace{3.5cm}} \\
\hline
$[\pi/2, 3\pi/4]$ & \underline{\hspace{3.5cm}} \\
\hline
$[3\pi/4, \pi]$ & \underline{\hspace{3.5cm}} \\
\hline
$[\pi, 5\pi/4]$ & \underline{\hspace{3.5cm}} \\
\hline
$[5\pi/4, 3\pi/2]$ & \underline{\hspace{3.5cm}} \\
\hline
$[3\pi/2, 7\pi/4]$ & \underline{\hspace{3.5cm}} \\
\hline
\end{tabular}

\vspace{6pt}
\textbf{Note:} $r = 1 + \cos\theta \geq 0$ for all $\theta$. What value of $\theta$ in
$[0, 2\pi]$ gives $r = 0$? \underline{\hspace{2.5cm}}

\end{notesbox}

\vspace{0.08in}

%----------------------------------------------------------------------
% Practice
%----------------------------------------------------------------------
\begin{practicebox}

\textbf{Guided Practice 1.}\quad The table below shows values of $r = 3\cos\theta$.

\vspace{4pt}
\renewcommand{\arraystretch}{1.5}
\begin{tabular}{|c|c|c|c|c|c|}
\hline
$\theta$ & $0$ & $\pi/4$ & $\pi/2$ & $3\pi/4$ & $\pi$ \\
\hline
$r$ & $3$ & $\approx 2.12$ & $0$ & $\approx -2.12$ & $-3$ \\
\hline
\end{tabular}

\vspace{6pt}
(a)\quad Classify each interval as P\&I, P\&D, N\&I, or N\&D.

\vspace{4pt}
$[0, \pi/4]$: \underline{\hspace{2.5cm}} \quad
$[\pi/4, \pi/2]$: \underline{\hspace{2.5cm}} \quad
$[\pi/2, 3\pi/4]$: \underline{\hspace{2.5cm}} \quad
$[3\pi/4, \pi]$: \underline{\hspace{2.5cm}}

\vspace{0.10in}
\textbf{Guided Practice 2.}\quad Using the table in Practice 1:

\vspace{4pt}
(a)\quad Compute the average rate of change of $r = 3\cos\theta$ from $\theta = 0$ to $\theta = \pi/2$.

\vspace{4pt}
ARC $= \dfrac{r(\pi/2) - r(0)}{\pi/2 - 0} = \dfrac{\rule{1.5cm}{0.4pt} - \rule{1.5cm}{0.4pt}}{\pi/2} = $ \blank{3cm}

\vspace{6pt}
(b)\quad Compute the average rate of change from $\theta = \pi/2$ to $\theta = \pi$.

\vspace{4pt}
ARC $= $ \blank{5cm}

\vspace{6pt}
(c)\quad On $[\pi/2, \pi]$, $r$ is negative. Does the curve move toward or away from the origin?
Explain using the N\&D or N\&I classification.

\writeline

\end{practicebox}

\end{document}
